\documentclass[11pt]{article}
\usepackage[margin=1in]{geometry}
\usepackage[T1]{fontenc}
\usepackage{lmodern,amsmath,amssymb,amsthm,booktabs,microtype,xcolor}
\usepackage[colorlinks=true,urlcolor=blue!50!black]{hyperref}
\newtheorem{theorem}{Theorem}
\hypersetup{pdftitle={How many perfect matchings test a ladder?}}
\title{\LARGE\bfseries How many perfect matchings test a ladder?}
\author{}
\date{Working note \(\,\cdot\,\) November 16, 2026}
\begin{document}
\maketitle
\noindent Choose a list of perfect matchings before seeing an edge coloring.
How short can the list be if it must detect a rainbow perfect matching
whenever one exists? For the ladder \(L_k=P_k\square K_2\) with \(k\ge5\)
columns, let \(h(k)\) denote the minimum list size when the coloring uses
exactly \(|E(L_k)|-3\) colors. A matching is rainbow if its edges have distinct
colors. We show that \(h(5)=7\) and \(h(k)=\Theta(\log k)\).

\section{Seven of the eight matchings}
Write \(r_i\) for the rung in column \(i\), numbering from zero, and \(H_i\)
for the pair of rails from column \(i\) to column \(i+1\), with upper rail
\(u_i\). Every perfect matching is determined by its set \(J\) of horizontal
gaps: select \(H_i\) for \(i\in J\) and fill uncovered columns with rungs.
The possible \(J\) are exactly the subsets with no consecutive indices,
as seen by examining the first free column.

\begin{theorem}On five columns the unique minimum test family consists of
all seven matchings with \(J\ne\varnothing\).\end{theorem}
\begin{proof}
The only matching omitted from this family uses all five rungs. Suppose it
is rainbow. Keep one edge of each color, choosing all five rungs among the
representatives. Exactly three edges are discarded. There are four disjoint
rail pairs, so some \(H_i\) is entirely kept. Replacing the rungs at
\(i,i+1\) by \(H_i\) produces a second rainbow matching.

For necessity, the table gives colorings having exactly one rainbow matching.
Each brace or rail pair in the right column is one color class. Every
unlisted edge has its own color. Each row has deficit three.
\begin{center}\small
\begin{tabular}{@{}ll@{}}
\toprule
Only rainbow gap set \(J\) & Nonsingleton classes\\
\midrule
\(\{0\}\) & \(\{r_0,r_4,u_2\},\ H_3\)\\
\(\{1\}\) & \(\{r_0,r_1,u_3\},\ H_0\)\\
\(\{0,2\}\) & \(\{r_0,r_2,r_4,u_3\}\)\\
\(\{0,3\}\) & \(\{r_0,r_1,u_1\},\ \{r_4,u_0\}\)\\
\bottomrule
\end{tabular}
\end{center}
To check each row, \(u_i\) occurs exactly when \(i\in J\), and \(r_i\)
occurs exactly when neither \(i-1\) nor \(i\) belongs to \(J\). Apply these
rules to the empty set, the four singletons, and the three pairs
\(\{0,2\},\{0,3\},\{1,3\}\). For example, in the first row gap 3 is
forbidden. Without it \(r_4\) occurs, so both \(r_0\) and \(u_2\) must be
absent. This forces gap 0 and excludes gaps 1 and 2. Reversing the first
three rows supplies the other three witnesses. Every successful list must
contain each of the seven matchings.
\end{proof}

\newpage
\section{Logarithmically many tests suffice and are necessary}
\begin{theorem}Let \(m=3k-2\) and \(k\ge5\). Then
\[
 \left\lceil\log_3\left\lfloor (k-1)/2\right\rfloor\right\rceil
 \ \le\ h(k)\ \le\
 \left\lfloor512\log\binom{m}{3}\right\rfloor+1 ,
\]
where the logarithm on the right is natural.\end{theorem}
\begin{proof}
For the lower bound, use the odd interior columns from 1 through \(k-2\).
In each test matching such a column has three possible states: a rung, a
horizontal tile extending left, or one extending right. A list of \(t\)
matchings gives each column a word of length \(t\) on these three states.

Suppose two chosen columns \(i<j\) have equal words. Give each of
\[
 \{r_i,r_j\},\qquad \{u_{i-1},u_{j-1}\},\qquad \{u_i,u_j\}
\]
its own repeated color, with fresh colors elsewhere. The spacing makes
the three pairs disjoint. A matching contains a repeated color precisely
when the two column states agree. Every test therefore fails, although
the single tile at gap \(i-1\) with rungs elsewhere is a rainbow matching.
Thus the \(\lfloor(k-1)/2\rfloor\) words must be distinct, proving the bound.

For the upper bound, start with a coloring and a rainbow perfect matching
\(M\). Keep one representative edge of every color, taking the edge in
\(M\) whenever possible. The discarded set \(F\) has exactly three edges,
and \(M\) avoids \(F\). Any matching avoiding \(F\) is rainbow.
It suffices to find one list that simultaneously avoids every avoidable
three-edge set in at least one of its tests.

Take independent fair bits on the gaps and set bits outside the ladder to
zero. Select precisely those gaps with bit one and both neighboring bits
zero. The resulting gaps are nonconsecutive and specify a perfect matching.
For each avoidable \(F\), fix one tiling avoiding it. For each forbidden
rung, that tiling covers its column with a horizontal tile. Prescribe a one
on that tile's gap and zeros on its neighbors. For each forbidden rail,
prescribe zero on its gap. All prescriptions agree with the fixed tiling's
gap word, including at the boundary. At most nine fair bits were prescribed.
With probability at least \(2^{-9}\) they all agree, producing a tiling
that avoids \(F\).

For \(t\) independent tests, a union bound over the avoidable triples gives
\[
 \Pr(\text{some avoidable triple is missed})\le
 \binom{m}{3}(1-2^{-9})^t\le\binom{m}{3}\exp(-t/512).
\]
For the stated upper bound this is less than one. A
successful list exists, and any duplicate tests may be removed.
\end{proof}

\small\noindent
\textbf{Method and verification.}
Prescribing a few binary coordinates is related to the universal sets of
Naor, Schulman and Srinivasan,
\href{https://authors.library.caltech.edu/records/wq91j-3w168}{\emph{Splitters and Near-optimal Derandomization}}.
The proof above gives its own probability estimate. Prior coverage of the
ladder bounds is unverified. The
\href{https://github.com/Akhil-rav/discrete-math-lab/blob/main/tests/test_ladder_test_families.py}{five-column check}
tests the witness table and every candidate list against all deficit-three
edge partitions.
\end{document}
